\chapter{Séries Temporais e Dinâmica de Painel Avançadas}

\section{Parâmetros variantes no tempo}

Em `tvp` o vetor de coeficientes é um passeio aleatório observado com
ruído:

\[
y_t = x_t'\beta_t + \varepsilon_t, \qquad
\beta_t = \beta_{t-1} + \eta_t
\]

`tvp` estima as variâncias de $\varepsilon_t$ e $\eta_t$ por MV e
reporta o caminho suavizado de $\beta_t$.

\begin{lstlisting}[caption={Regressão com parâmetros variantes no tempo}]
let n = 200
generate df x = rnormal(n)
generate df y = (1 + 0.02*_n) + (0.5 - 0.003*_n)*x + rnormal(n)

let m_tvp = tvp(y ~ x, df)
print(m_tvp)
\end{lstlisting}

Use `tvp_var` para uma VAR completa com coeficientes variantes no tempo.

\section{VAR de limiar e de mudança de regime Markov}

`setar` permite dois regimes AR divididos por um valor de limiar da
variável dependente defasada.

\begin{lstlisting}[caption={SETAR}]
let m_setar = setar(y, df, order=2, delay=1)
print(m_setar)
\end{lstlisting}

`msvar` permite que os coeficientes mudem de acordo com uma cadeia
latente de Markov.

\begin{lstlisting}[caption={VAR com mudança de regime Markov}]
let m_msvar = msvar(y1 + y2, df, regimes=2, lags=1)
print(m_msvar)
\end{lstlisting}

\section{Dinâmicas assimétricas e não lineares}

`nardl` decompõe um regressor em somas parciais positivas e negativas,
permitindo respostas assimétricas de longo prazo.

\begin{lstlisting}[caption={NARDL}]
let m_nardl = nardl(y ~ x, df, lags=1)
print(m_nardl)
\end{lstlisting}

`pstr` faz o mesmo em painel, com coeficientes que mudam suavemente de
acordo com uma função de transição logística de uma variável observada.

\begin{lstlisting}[caption={Regressão de Transição Suave em Painel}]
let m_pstr = pstr(y ~ x1 + x2, df, q="z", id="firm")
print(m_pstr)
\end{lstlisting}

\section{Coeficientes funcionais}

`fcoef` estima coeficientes que variam suavemente com uma variável
moderadora $z$, usando regressão kernel local linear.

\begin{lstlisting}[caption={Modelo de coeficientes funcionais}]
let m_fcoef = fcoef(y ~ x1 + x2, df, z="z", points=20)
print(m_fcoef)
\end{lstlisting}

\section{Bayesian VAR e VAR de frequência mista}

`bvar` estima uma VAR com priori de Minnesota, encolhendo defasagens
próprias em direção a 1 e defasagens cruzadas em direção a 0.

\begin{lstlisting}[caption={Bayesian VAR}]
let m_bvar = bvar(y1 ~ y2, df, lags=1)
print(m_bvar)
\end{lstlisting}

`mfvar` combina variáveis de baixa e alta frequência usando agregação
Almon MIDAS exponencial.

\begin{lstlisting}[caption={VAR de frequência mista}]
let m_mfvar = mfvar(df_q, yq, df_m, ym, agg=3, lags=1)
print(m_mfvar)
\end{lstlisting}

\section{VAR em painel}

`pvar` estima uma VAR em painel usando GMM estilo Arellano-Bond:
valores defasados em níveis instrumentam as equações em primeiras
diferenças.

\begin{lstlisting}[caption={VAR em painel}]
let m_pvar = pvar(y1 ~ y2, df, id="firm", lags=1)
print(m_pvar)
\end{lstlisting}

Esses comandos são poderosos, mas sensíveis a convenções. Sempre
compare a saída com uma implementação de referência quando possível e
inspecione as probabilidades de transição, autovalores ou estatísticas
J reportadas pelo estimador.
